%% GENERATED by tools/from_docs.py from stepss-docs. Do not edit by hand:
%% edit the page named below in stepss-docs and regenerate.
%% source: stepss-docs/src/content/docs/getting-started/quickstart.mdx

\chapter{Three ways to drive the engines}\label{doc:quickstart}

You drive STEPSS one of two ways, and both are covered here.

{\def\LTcaptype{none} % do not increment counter
\begin{small}\begin{longtable}[]{@{}
  >{\raggedright\arraybackslash}p{(\linewidth - 8\tabcolsep) * \real{0.1429}}
  >{\centering\arraybackslash}p{(\linewidth - 8\tabcolsep) * \real{0.2381}}
  >{\centering\arraybackslash}p{(\linewidth - 8\tabcolsep) * \real{0.2381}}
  >{\centering\arraybackslash}p{(\linewidth - 8\tabcolsep) * \real{0.2381}}
  >{\raggedright\arraybackslash}p{(\linewidth - 8\tabcolsep) * \real{0.1429}}@{}}
\toprule\noalign{}
\begin{minipage}[b]{\linewidth}\raggedright
Edition
\end{minipage} & \begin{minipage}[b]{\linewidth}\centering
Dynamic Simulator (RAMSES)
\end{minipage} & \begin{minipage}[b]{\linewidth}\centering
Power Flow (Helios)
\end{minipage} & \begin{minipage}[b]{\linewidth}\centering
CODEGEN
\end{minipage} & \begin{minipage}[b]{\linewidth}\raggedright
Use it for
\end{minipage} \\
\midrule\noalign{}
\endhead
\bottomrule\noalign{}
\endlastfoot
\textbf{STEPSS GUI} & yes & yes & yes & Interactive work: load a case, run it, watch curves, build models \\
\textbf{STEPSS in Python} & yes & yes & & Scripting, parameter sweeps, and the scientific Python stack \\
\end{longtable}\end{small}
}

Those two are what you install; see
\hyperref[ch-install]{Installation}. Pick the GUI if you want to see
the system respond, Python if you want to automate. Neither wraps the other, so a
case built in one runs unchanged in the other.

New to STEPSS? \hyperref[doc:gui-first-run]{First Run} opens a bundled test system and
simulates it without your having to prepare any data.

The engines can also be invoked directly from a shell, which is what the
\textbf{Command line} tab on each section below documents. It is not a third thing to
install and most readers will not need it: the executables are not published
separately, they ship inside STEPSS GUI, which unpacks them while it runs. It
matters when you are wiring STEPSS into a job scheduler or another program.

\begin{center}\rule{0.5\linewidth}{0.5pt}\end{center}

\section{Dynamic Simulation (RAMSES)}\label{quickstart:dynamic-simulation-ramses}

\textbf{STEPSS GUI}

The window is six tabs, used left to right:

\begin{enumerate}
\def\labelenumi{\arabic{enumi}.}
\item
  \textbf{Launch STEPSS}, from your applications if you used an installer, or by double-clicking \texttt{stepss.jar}
\item
  \textbf{Get a case.} \textbf{File}, then \textbf{Open Examples} ships three complete test
  systems and fills in every file slot. Otherwise load your own on the \textbf{System
  Data} tab, which takes the system data files and the disturbance file
\item
  On \textbf{Observables}, tick what to record, and optionally set a quantity to plot
  live while the run proceeds
\item
  On \textbf{Power Flow Simulation}, \textbf{Run power flow} to get the operating point
\item
  On \textbf{Dynamic Simulation}, \textbf{Run dynamic simulation}
\item
  On \textbf{Analysis}, \textbf{Extract curves} to plot the result
\end{enumerate}

The GUI coordinates the power flow and RAMSES through lock files (\texttt{.lock\_RAMSES}, \texttt{.kill\_RAMSES}) for graceful process management.

For the full walkthrough on a real case, see
\hyperref[doc:gui-running]{Running a Simulation}.

\textbf{STEPSS in Python}

\begin{enumerate}
\def\labelenumi{\arabic{enumi}.}
\item
  \textbf{Install stepss}:

\begin{Verbatim}
pip install stepss
\end{Verbatim}
\item
  \textbf{Run a simulation}:

\begin{Verbatim}
import stepss

# Configure the case
case = stepss.cfg()
case.addData("data.dat")
case.addData("volt_rat.dat")
case.addData("settings.dat")
case.addObs("obs.obs")
case.addDst("disturbance.dst")

# Run
ram = stepss.sim()
ram.execSim(case, 150.0)
\end{Verbatim}
\item
  \textbf{Extract and plot results}:

\begin{Verbatim}
ext = stepss.extractor(case.getTrj())

# Plot bus voltage magnitude
ext.getBus('1041').mag.plot()

# Plot synchronous machine speed
ext.getSync('g1').S.plot()
\end{Verbatim}
\end{enumerate}

See the \hyperref[doc:py-api]{Python API Reference} for the complete interface.

\textbf{Command Line}

\textbf{The RAMSES command-line interface}

The RAMSES executable accepts the following arguments:

\begin{Verbatim}
ramses.exe [-t command_file] [-v]
\end{Verbatim}

{\def\LTcaptype{none} % do not increment counter
\begin{small}\begin{longtable}[]{@{}
  >{\raggedright\arraybackslash}p{(\linewidth - 2\tabcolsep) * \real{0.3158}}
  >{\raggedright\arraybackslash}p{(\linewidth - 2\tabcolsep) * \real{0.6842}}@{}}
\toprule\noalign{}
\begin{minipage}[b]{\linewidth}\raggedright
Flag
\end{minipage} & \begin{minipage}[b]{\linewidth}\raggedright
Description
\end{minipage} \\
\midrule\noalign{}
\endhead
\bottomrule\noalign{}
\endlastfoot
\texttt{-t\ command\_file} & Read inputs from a command file (non-interactive mode) \\
\texttt{-v} & Print version information and exit \\
\end{longtable}\end{small}
}

\textbf{Interactive mode} (no \texttt{-t} flag): RAMSES prompts for each input sequentially:

\begin{Verbatim}
Dump settings, comments and initialization data in file [press enter to skip] :
Disturbance file name [compulsory] :
Dump system trajectory in file [press enter to skip] :
Observable file name [compulsory] :
Dump continuous simulation trace in file [press enter to skip] :
Dump discrete simulation trace in file [press enter to skip] :
\end{Verbatim}

The observables prompt is shown only when a trajectory file was given; it is
labelled compulsory because at that point it is.

\textbf{Non-interactive mode} (\texttt{-t} flag): The command file provides the same inputs as lines, in order:

\begin{enumerate}
\def\labelenumi{\arabic{enumi}.}
\item
  \textbf{Prepare data files}: Network data, dynamic data, solver settings, LFRESV (power flow solution)
\item
  \textbf{Create a command file} (\texttt{cmd.txt}):

\begin{Verbatim}
data.dat volt_rat.dat settings.dat

disturbance.dst
result.rtrj
obs.obs
cont.trace
disc.trace
\end{Verbatim}

  The command file contains, in order:

  \begin{itemize}
  \tightlist
  \item
    \textbf{Data files} (one or more, space- or newline-separated), read until a blank line
  \item
    \textbf{Initialization dump file} (blank line to skip)
  \item
    \textbf{Disturbance file} (required)
  \item
    \textbf{Trajectory file} (blank line to skip)
  \item
    \textbf{Observables file} (required if trajectory file is given)
  \item
    \textbf{Continuous trace file} (blank line to skip)
  \item
    \textbf{Discrete trace file} (blank line to skip)
  \end{itemize}
\item
  \textbf{Run the simulation}:

\begin{Verbatim}
ramses.exe -tcmd.txt
\end{Verbatim}
\item
  \textbf{Process results}: Use the stepss extractor, DYNGRAPH, or other tools to read the trajectory file
\end{enumerate}

\begin{docsnote}{Note}

Lines starting with \texttt{\#} in the command file are treated as comments and skipped.

\end{docsnote}

\begin{center}\rule{0.5\linewidth}{0.5pt}\end{center}

\section{Power Flow (Helios)}\label{quickstart:power-flow-helios}

\textbf{STEPSS GUI}

\begin{enumerate}
\def\labelenumi{\arabic{enumi}.}
\item
  \textbf{Launch STEPSS} and load the network and power-flow data files on the
  \textbf{System Data} tab
\item
  Go to \textbf{Power Flow Simulation} and click \textbf{Run power flow}. The pane reports bus
  voltages and angles, the generator table with its limits, and the system power
  balance
\item
  Inspect the result through \textbf{Bus overview}, \textbf{Branch flows}, \textbf{Generators \&
  SVCs}, \textbf{Adjustable transformers} and \textbf{Global power balance}, which enable
  once the run succeeds
\item
  \textbf{Add Helios results to data} feeds the solution back into the loaded case,
  ready for the dynamic simulation
\end{enumerate}

See \hyperref[doc:gui-interface]{The Interface} for what each control does.

\textbf{Command Line}

\textbf{The Helios command-line interface}

The Helios executable accepts the following arguments:

\begin{Verbatim}
helios [-t command_file] [--plain] [--tui]
\end{Verbatim}

{\def\LTcaptype{none} % do not increment counter
\begin{small}\begin{longtable}[]{@{}
  >{\raggedright\arraybackslash}p{(\linewidth - 2\tabcolsep) * \real{0.3158}}
  >{\raggedright\arraybackslash}p{(\linewidth - 2\tabcolsep) * \real{0.6842}}@{}}
\toprule\noalign{}
\begin{minipage}[b]{\linewidth}\raggedright
Flag
\end{minipage} & \begin{minipage}[b]{\linewidth}\raggedright
Description
\end{minipage} \\
\midrule\noalign{}
\endhead
\bottomrule\noalign{}
\endlastfoot
\texttt{-t\ command\_file} & Read inputs from a command file (non-interactive mode) \\
\texttt{-\/-plain} & Force plain-text output instead of the interactive terminal UI \\
\texttt{-\/-tui} & Force the interactive terminal UI \\
\end{longtable}\end{small}
}

\textbf{Interactive mode} (no \texttt{-t} flag): Helios prompts for data files, then presents a command menu. On a TTY it uses a colour-coded terminal UI; when stdin is a pipe it falls back to plain text automatically.

\textbf{Non-interactive mode} (\texttt{-t} flag): The command file provides data file names (one per line, blank line to end), followed by menu commands. Results go to stdout and errors to stderr with line numbers. Paths resolve relative to the working directory, so \texttt{cd} into the data directory first.

\begin{enumerate}
\def\labelenumi{\arabic{enumi}.}
\item
  \textbf{Prepare data files}: Network data and power flow data (generators, loads, SVC, LTC-V, PSHIFT-P)
\item
  \textbf{Create a command file} (\texttt{pf\_cmd.txt}):

\begin{Verbatim}
network.dat
powerflow.dat

VT
volt_rat.dat
E
\end{Verbatim}
\item
  \textbf{Run the power flow}:

\begin{Verbatim}
helios -t pf_cmd.txt > output.txt
\end{Verbatim}
\end{enumerate}

The run's outcome is reported through the process exit status: \texttt{0} converged, \texttt{1} input or usage error, \texttt{2} solve did not converge. See \hyperref[ch-pf]{Exit Status}.

Helios then presents an interactive menu. \texttt{VT} is the command that writes the
LFRESV file needed to initialize RAMSES; see
\hyperref[ch-pf]{Interactive Menu Commands}
for the full list.

\begin{center}\rule{0.5\linewidth}{0.5pt}\end{center}

\section{Typical Workflow}\label{quickstart:typical-workflow}

\begin{enumerate}
\def\labelenumi{\arabic{enumi}.}
\tightlist
\item
  \textbf{Define the network}: Specify buses, lines, transformers, and shunts
\item
  \textbf{Set up power flow data}: Define generators, loads, and the slack bus
\item
  \textbf{Run the power flow}: Compute the initial operating point and export LFRESV
\item
  \textbf{Add dynamic models}: Specify synchronous machines, excitation systems, speed governors, loads, and controllers
\item
  \textbf{Configure solver}: Set integration method, time steps, and tolerances
\item
  \textbf{Define disturbances}: Specify faults, line trips, parameter changes, etc.
\item
  \textbf{Run RAMSES}: Execute the dynamic simulation
\item
  \textbf{Analyze results}: Extract and visualize trajectories of voltages, frequencies, and powers
\end{enumerate}

\section{Next Steps}\label{quickstart:next-steps}

\begin{itemize}
\tightlist
\item
  \hyperref[doc:gui-first-run]{First Run}, Open a bundled test system in STEPSS GUI
\item
  \hyperref[doc:gui-running]{Running a Simulation}, A complete run, start to finish
\item
  \hyperref[doc:py-examples]{Python Examples}, The same ground in a script
\item
  \hyperref[ch-files]{File Formats}, Learn about data file syntax
\item
  \hyperref[ch-network]{Network Modeling}, Define your power system network
\item
  \hyperref[ch-pf]{Power Flow}, Power flow data and control parameters
\item
  \hyperref[ch-sync]{Dynamic Data Records}, Add generators, loads, and controllers
\end{itemize}
